% Responsable: P4
% Tablas y figuras: python -m evaluation.metrics  (results/tablas/*.tex, results/figures/*.png).
% \tablagenerada y \figuragenerada están definidas en informe.tex (copias en tablas/ y figuras/).
\section{Resultados}
\label{sec:resultados}


Se comparan cinco sistemas sobre las mismas 73 oraciones: nuestro parser, spaCy, Transformers
(XLM-R para las etiquetas y BETO para decidir la gramaticalidad, sección~\ref{sec:comparacion_nlp})
y el LLM en sus dos versiones, \emph{zero-shot} y \emph{few-shot}
(sección~\ref{sec:analisis_llm}). Cada sistema aparece solo en las tablas para las que produce
datos: spaCy no decide si una oración es válida, y Transformers no da funciones sintácticas ni
unión de la frase preposicional. Esta sección presenta las cifras; las secciones
\ref{sec:comparacion_nlp} a~\ref{sec:comparacion_enfoques} las interpretan.

\subsection{Aceptación}

\begin{table}[h]
  \centering
  \small
  \tablagenerada{aceptacion}\par\medskip
  \tablagenerada{aceptacion_espanol}
  \caption{Aceptación contra \texttt{en\_gramatica} (arriba) y contra \texttt{espanol\_valido}
    (abajo). ``---'': la métrica no aplica al sistema.}
  \label{tab:aceptacion}
\end{table}

\begin{figure}[h]
  \centering
  \figuragenerada{aceptacion}
  \caption{Métricas de aceptación contra las reglas de la gramática.}
  \label{fig:aceptacion}
\end{figure}

El parser propio acepta exactamente las 45 oraciones de la gramática y rechaza las 28 que no
lo están. Además, todas las rechaza por la razón esperada: 15 por error sintáctico y 13 por
concordancia (tabla~\ref{tab:tipo_error}). Este resultado no mide qué tan bueno es el enfoque,
porque el conjunto se etiquetó con la misma gramática que implementa el parser. Lo que muestra
es que la implementación es fiel a la especificación, incluso en los casos difíciles: la
concordancia a distancia (\emph{El perro de los niños come} se acepta y \emph{*...comen} se
rechaza) y la palabra ambigua \emph{bajo}.

La comparación interesante es contra \texttt{espanol\_valido}. Ahí el parser tiene precisión
1.00, pero su cobertura baja a 0.80. Nunca acepta español incorrecto, pero rechaza 11 de las 56
oraciones correctas: todo lo que la CFG no genera (sujeto omitido, orden libre, coordinación,
subordinadas y palabras fuera del léxico). Ese 20\,\% es la medida concreta de cuánto español
correcto queda fuera de la gramática.

\begin{figure}[h]
  \centering
  \figuragenerada[0.7\linewidth]{reglas_vs_espanol}
  \caption{Exactitud de la aceptación contra las dos referencias.}
  \label{fig:reglas_vs_espanol}
\end{figure}

Los otros dos sistemas que deciden se comportan de forma opuesta. Transformers tiene
exactitud 0.85 contra el español y solo 0.73 contra las reglas: acepta 10 de las 11 oraciones
correctas que la gramática no genera, es decir, sigue su conocimiento del español. El LLM hace
lo contrario: 0.97 (\emph{zero-shot}) y 0.99 (\emph{few-shot}) contra las reglas, y 0.82 y 0.86
contra el español. Cuando recibe la gramática en el prompt, rechaza casi todo lo que la gramática
no genera, aunque sea español correcto. Es el sistema más cercano al parser en la decisión, y el
único además del parser que dice por qué rechaza: acierta el tipo de error en el 96\,\% y el
100\,\% de los rechazos.

\begin{table}[h]
  \centering
  \small
  \resizebox{\linewidth}{!}{\tablagenerada{tipo_error}}
  \caption{Tipo de error esperado frente a la respuesta del sistema, en las 28 oraciones fuera
    de la gramática.}
  \label{tab:tipo_error}
\end{table}

\subsection{Resultados por nivel}

\begin{table}[h]
  \centering
  \small
  \tablagenerada{exactitud_aceptacion_por_nivel}\par\medskip
  \tablagenerada{etiquetado_por_nivel}
  \caption{Exactitud de la aceptación (arriba) y del etiquetado (abajo) por nivel.}
  \label{tab:por_nivel}
\end{table}

\begin{figure}[h]
  \centering
  \figuragenerada{por_nivel}
  \caption{Exactitud por nivel de complejidad.}
  \label{fig:por_nivel}
\end{figure}

Los niveles 1 y 2 son fáciles para todos los sistemas que deciden, salvo Transformers, que ya
en el nivel 1 rechaza una de cada cuatro oraciones correctas (exactitud 0.75): BETO castiga las
oraciones cortas y con palabras poco frecuentes (sección~\ref{sec:comparacion_nlp}). En el
nivel 3 las diferencias se hacen grandes. El parser sigue en 1.00 contra las reglas; el LLM baja
a 0.96 en las dos versiones, por una sola oración en cada una (\emph{zero-shot} rechaza \emph{Él
mató un elefante en sus pijamas}; \emph{few-shot} acepta \emph{Leímos un libro}); y
Transformers cae a 0.54, porque es el nivel con las 11 oraciones correctas fuera de la
gramática, que acepta casi todas. En el etiquetado el patrón se repite con menos fuerza: spaCy es
el único que pierde exactitud desde el nivel 1 (0.92), por las oraciones cortas, y los demás
solo bajan en el nivel 3.

\subsection{Etiquetado y funciones sintácticas}

\begin{table}[h]
  \centering
  \small
  \tablagenerada{etiquetado}\quad\tablagenerada{funciones}
  \caption{Exactitud del etiquetado (y oraciones con otro número de palabras), y de las
    funciones sintácticas.}
  \label{tab:etiquetado}
\end{table}

\begin{table}[h]
  \centering
  \small
  \tablagenerada{etiquetado_categoria}
  \caption{Cobertura (\emph{recall}) del etiquetado por categoría.}
  \label{tab:etiquetado_categoria}
\end{table}

\begin{figure}[h]
  \centering
  \figuragenerada[0.8\linewidth]{etiquetado_categoria}
  \caption{Cobertura del etiquetado por categoría.}
  \label{fig:etiquetado_categoria}
\end{figure}

El parser etiqueta sin errores, y da bien el sujeto, el verbo y el objeto, en todas las
oraciones que puede derivar (58 de 73: las aceptadas y las que fallan solo en la
concordancia). En las otras 15 no produce etiquetas. Los demás sistemas etiquetan las 73. El LLM
es el más exacto (0.99 en las dos versiones), seguido de Transformers (0.98) y de spaCy (0.94),
cuyo punto débil son los verbos: su cobertura en \textit{Verb} es 0.82, porque en oraciones
cortas como \emph{El perro come} etiqueta el verbo como sustantivo. En la categoría \texttt{X}
(palabras fuera del léxico) spaCy y Transformers tienen cobertura 0.50, pero no por error: dan la
categoría real de palabras como \emph{que} o \emph{no}, que el conjunto marca como \texttt{X}
por convención. El LLM, que tiene el léxico en el prompt, sí las marca \texttt{X}.

En las funciones sintácticas el LLM \emph{few-shot} queda cerca del parser (0.97, 0.96 y 1.00
en sujeto, verbo y objeto) y por encima de spaCy (0.88, 0.86 y 0.92). La exactitud del parser
es 1.00 pero, otra vez, solo sobre las oraciones que puede derivar; la
sección~\ref{sec:comparacion_nlp} compara los sistemas en el mismo subconjunto.

\subsection{Unión de la frase preposicional}

\begin{table}[h]
  \centering
  \small
  \tablagenerada{pp_adjunta}\par\medskip
  \tablagenerada{pp_confusion}
  \caption{Unión de la primera frase preposicional en las 17 oraciones que la tienen.}
  \label{tab:pp}
\end{table}

\begin{figure}[h]
  \centering
  \figuragenerada{funciones_pp}
  \caption{Funciones sintácticas y unión de la frase preposicional.}
  \label{fig:funciones_pp}
\end{figure}

Aquí aparece el costo del conflicto LL(1) de la sección~\ref{sec:gramatica}. El parser une la
frase preposicional al sustantivo más cercano siempre que hay un objeto. Por eso acierta en
las 6 oraciones donde la unión correcta es al sustantivo, y falla en las 4 donde un hablante la
une al verbo (\emph{Ellos ven los árboles altos desde la casa}). Las 2 oraciones ambiguas las
resuelve sin avisar, como si no lo fueran, así que detecta 0 de 2 ambigüedades. Su exactitud es
0.59 (10 de 17). Una de las fallas es \emph{El perro corre en el parque}, que no se puede
derivar porque \emph{corre} no está en el léxico.

Los sistemas estadísticos aciertan más porque deciden por el sentido y no por la posición. spaCy
acierta 13 de 17: falla en dos oraciones donde une la frase al sustantivo cuando un hablante la
une al verbo (\emph{Ellos ven los árboles altos desde la casa}). El LLM es el mejor, con 15 de 17
en \emph{zero-shot} y 14 de 17 en \emph{few-shot}: une \emph{con su perro} al verbo y \emph{de
los niños} al sustantivo, como un hablante. Pero ningún sistema detecta las 2 oraciones
ambiguas. spaCy y el LLM eligen siempre el verbo, y el LLM nunca responde \texttt{ambigua},
aunque el prompt le ofrece esa opción.
